% Capítulo 5, Considerações Finais
% "Considerações Finais" (não "Conclusão") pela norma da biblioteca.

% \section{Síntese do trabalho até a etapa atual}
\label{sec-sintese}

O trabalho avançou ao longo do Trabalho de Conclusão de Curso~I em três linhas de trabalho complementares. A primeira foi a delimitação do problema, com a formulação da questão central, dos invariantes alvo e dos critérios de avaliação a partir da literatura normativa e da Proposta original do trabalho. A segunda foi o mapeamento sistemático da literatura, que produziu um protocolo replicável, executou a busca em uma plataforma agregadora e gerou um conjunto de 165 estudos primários para leitura completa, com o esquema de classificação descrito na Seção~\ref{sec-ms}. A terceira foi a construção dos artefatos da vertente experimental e a execução da verificação formal: o repositório \texttt{tcc-desenvolvimento} versiona o \emph{cluster} Infinispan parametrizado, o programa de carga em Java sobre Hot Rod, a especificação TLA+ \texttt{SessionStore.tla} com os invariantes I1 a I6 e o \emph{pipeline} de análise estatística em Python, todos detalhados no Capítulo~\ref{cap-desenvolvimento}. A verificação formal foi executada em cinco rodadas do TLC, com contraexemplos produzidos para I1, I2, I5 e I6 e ausência de violação para I3 e I4 nas configurações exploradas, conforme reportado na Seção~\ref{sec-mc-resultados}. Os cenários de \revini{CON-14}parada ordenada de nó (F1, por \texttt{podman stop} e \texttt{podman start})\revfim{CON-14} e de atraso com \emph{jitter} de canal (F3) do \emph{baseline} experimental foram executados contra o \emph{cluster} vivo nesta etapa, em máquina virtual Linux\revini{CON-01}. Essa primeira bateria de medição, de junho, não mediu o ciclo de vida de sessões, por defeitos da carga e da análise identificados em outubro (Seção~\ref{sec-baseline-estado}).\revfim{CON-01}

\revini{CON-02}Na versão parcial do Trabalho de Conclusão de Curso~II, o trabalho chega com a vertente formal em estágio mais avançado que a experimental. \revini{CON-23}A Seção~\ref{sec-revisoes} reúne as mudanças em relação ao plano original e seus motivos.\revfim{CON-23} Na vertente formal, o modelo TLA+ do protocolo de sessão inclui quedas, recuperações e partições. Na verificação pelo TLC, as fórmulas globais de I1 e I2 são violadas por construção em um modelo com várias réplicas, o que levou a propriedades avaliadas réplica a réplica. No cenário semeado, a irreversibilidade da invalidação por réplica só foi violada sob partição, coerente com a perda, na fusão, de escritas feitas na partição menor. Com \emph{login} habilitado, ela também é violada sem falha, por artefatos do modelo. As formas de I3 e I4 declaradas no modelo, de faixa e de tipo, não foram violadas, com verificação exaustiva apenas na menor configuração\revini{CON-29}\revini{CON-34}, \texttt{tiny}, acrescentada nas verificações de setembro de 2026, depois das cinco rodadas do Trabalho de Conclusão de Curso~I\revfim{CON-34}\revfim{CON-29}, e a imutabilidade do motivo, avaliada réplica a réplica, é violada. I5 e I6 são violados por dois mecanismos: o empate de marcas de tempo, defeito do modelo cuja correção está pendente, e, sob partição, a perda, na fusão, de uma escrita ainda não propagada, com divergência permanente (Seção~\ref{sec-mc-resultados}).

Na vertente experimental, o ambiente está montado e documentado (Seção~\ref{sec-arquitetura}): o \emph{cluster} de três nós, o programa de carga com os cenários S1, S2 e S3, os injetores das quatro classes de falha e os \emph{scripts} de análise. Os defeitos da primeira bateria foram corrigidos, e uma nova medição de referência, com os modos síncrono e assíncrono, foi pré-registrada, mas ainda não executada. Todas as execuções na máquina virtual até aqui ocorreram em \texttt{DIST\_SYNC}, de modo que o comportamento do modo assíncrono, alvo do trabalho, ainda não foi medido. O injetor de partição foi verificado sem carga de aplicação, e o de degradação parcial passou apenas pelas verificações que não afetam o tráfego, sem injeção real até aqui.

O controlador heurístico está projetado e implementado em parte (Seção~\ref{sec-controlador}). O gatilho primário combina latência de cauda relativa, taxa de erro e topologia, e as violações de I1, I2 e I4 formam o gatilho secundário. A convergência (I5 e I6) entra apenas como desfecho. A primeira etapa atua apenas sobre parâmetros do cliente, e a atuação sobre a configuração do \emph{cluster} é uma segunda etapa planejada, cujos parâmetros candidatos ainda não foram verificados. O controlador ainda não foi ligado ao laço de carga, e sua avaliação, desenhada em blocos pareados (Seção~\ref{sec-avaliacao-controlador}), depende da nova medição de referência e da campanha com S3.\revfim{CON-02}

O conjunto de invariantes alvo, irreversibilidade de invalidação por \emph{logout}, consistência de bloqueio, monotonicidade local do contador, unicidade do motivo de invalidação, convergência do estado lógico e convergência do contador, encontra respaldo direto em \revini{CON-03}\citet{nist_sp_800_63b4_2025}\revfim{CON-03} e em \citet{owasp_session_mgmt_cheatsheet} para as propriedades de sessão e na literatura formal sobre replicação otimista, em particular \citet{shapiro_crdt_2011} e \citet{nair_safety_distributed_2020}, para as propriedades de convergência. A escolha do par de métricas latência de cauda e taxa de violação como contribuição principal é sustentada pela síntese inicial do mapeamento sistemático: a primeira métrica é o padrão consolidado da literatura para sistemas com cauda longa, e a segunda métrica não aparece como métrica primária de avaliação experimental \revini{CON-17}nos cerca de trinta trabalhos lidos na íntegra até esta etapa, constatação que a classificação dos 165 estudos do mapeamento ainda precisa confirmar\revfim{CON-17}. \revini{CON-04}O ponto comparativo mais próximo, \citet{nguyen_weaker_consistency_2019}, combina consistência relaxada com monitoramento e recuperação reativos, sem verificação formal a montante.\revfim{CON-04} Os contraexemplos do TLC fornecem o catálogo inicial de padrões de divergência a serem mitigados pelo controlador heurístico previsto para a etapa seguinte. \revini{CON-05}Até a correção do empate de marcas de tempo, os contraexemplos de I5 e I6 que não passam pela fusão não informam sobre o protocolo.\revfim{CON-05}

%\section{Limitações conhecidas}
%\label{sec-limitacoes}

Três limitações merecem registro nesta etapa. A primeira é a execução do mapeamento sistemático restrita a uma plataforma agregadora, com a complementação por busca manual nas bases originais formalizada como ameaça à validade na Seção~\ref{sec-ms} e integrada às atividades remanescentes para o Trabalho de Conclusão de Curso~II.
\revini{CON-15}A segunda é a leitura na íntegra das referências\revfim{CON-15} \revini{CON-18}citadas, que não alcançou todas. Por restrição de acesso no período, oito das fontes citadas foram consultadas apenas pelo resumo ou por trechos: \citet{newcombe_aws_tla_2014}, \citet{liu_quantitative_consistency_2015}, \citet{naskos_probabilistic_mc_2017}, \citet{wu_anna_autoscale_2020}, \citet{gambi_assurance_selfadaptive_2013}, \citet{rosa_selfmgmt_2013}, \citet{enea_must_2024} e \citet{benmokhtar_scheduling_2024}; as afirmações que dependem delas se limitam ao que os resumos sustentam. A resolução final dessas exceções compõe\revfim{CON-18} a curadoria bibliográfica do Trabalho de Conclusão de Curso~II. A terceira é a execução parcial dos artefatos. A vertente formal completou as cinco rodadas planejadas do TLC nas três configurações da Tabela~\ref{tab-parametros-baseline} (T21)\revini{CON-28}\revini{CON-35}; a configuração \texttt{tiny} foi acrescentada depois, nas verificações de setembro de 2026 (Seção~\ref{sec-mc-resultados})\revfim{CON-35}\revfim{CON-28}. \revini{CON-32}A configuração do \emph{cluster}, o programa de carga em Java, os \emph{scripts} de injeção de falhas e o \emph{pipeline} de análise estatística estão implementados. O coletor de métricas OpenMetrics \citep{openmetrics_spec} foi escrito, mas não foi integrado ao código principal.\revfim{CON-32}\revini{CON-16}\revfim{CON-16} O \emph{baseline} experimental teve os cenários de \revini{CON-21}parada ordenada de nó (F1)\revfim{CON-21} e de atraso com \emph{jitter} de canal (F3) executados contra o \emph{cluster} vivo, em máquina virtual Linux com o módulo \texttt{sch\_netem} disponível, \revini{CON-06}\revfim{CON-06}na forma de um piloto reduzido em relação à bateria plena dos parâmetros T9 e T10, que segue como protocolo a escalar. \revini{CON-07}Essa bateria não mediu o ciclo de vida de sessões, e nenhum número de latência de junho é usado como resultado (Seção~\ref{sec-baseline-estado}).\revfim{CON-07}

\revini{CON-30}No Trabalho de Conclusão de Curso~I, a abstração do canal de replicação, no modelo formal, capturava atraso e reordenação mas não modelava perdas permanentes ou particionamento. Essa simplificação era proposital naquela etapa, dado que a verificação visava o fluxo nominal sob entrelaçamento concorrente; o refinamento com falhas explícitas integrava o plano para o Trabalho de Conclusão de Curso~II (Apêndice~\ref{ap-plano-tcc2}, etapas 3 e 4).\revfim{CON-30} \revini{CON-08}No Trabalho de Conclusão de Curso~II, o modelo ganhou as ações de queda, recuperação, partição e fusão (Seção~\revini{CON-33}\ref{sec-spec-formal}\revfim{CON-33}). O canal do modelo não perde mensagens na versão integrada, e a variante com perda tem caráter ilustrativo.\revfim{CON-08}

\section*{\revini{CON-09}Limitações e ameaças à validade}
\label{sec-limitacoes}

Na vertente formal, os resultados de I3 e I4 nas configurações maiores valem só até o tempo-limite das execuções, e as cardinalidades verificadas são pequenas, recorte apoiado pela observação de que \revini{CON-20}98\,\% das falhas de produção analisadas\revfim{CON-20} se manifestam com três nós ou menos \citep{yuan_simple_testing_2014}, mas que não exclui defeitos em escalas maiores. Os contraexemplos de I5 e I6 que não passam pela fusão decorrem do empate de marcas de tempo e não terão interpretação enquanto ele não for corrigido.

Na vertente experimental, nenhum dado de latência ou de violação é ainda utilizável: a bateria de junho não criou sessões, e a nova medição não foi executada. Os F1 já executados foram paradas ordenadas, e não quedas silenciosas, e o modo silencioso não foi ensaiado na máquina virtual. A partição foi exercida apenas sem carga de aplicação, e o tempo de recomposição da segunda repetição, de 188\,s, obtido em condições diferentes das demais (\emph{cluster} com histórico e coleta de monitoramento intensa), tem causa não isolada entre duas hipóteses registradas. A nova medição usa dez repetições por célula, e não as trinta do plano, o que reduz o poder estatístico. Ela mantém ainda 64 segmentos, sem fonte que justifique o valor, e usa distribuição normal em F3, no lugar da lognormal. Os parâmetros do cenário S3 são premissas de modelagem, e a estimativa de energia depende de um modelo de potência e do tempo de CPU virtual, com as limitações da Seção~\ref{sec-energia}.

O ambiente de execução é uma ameaça à validade de outra natureza. O computador hospedeiro registrou dezessete desligamentos não ordenados desde setembro de 2026, dois deles com parada do sistema por erro registrada. O autor atribui os desligamentos a quedas no fornecimento de energia elétrica na região. Uma dessas quedas, em 18 de outubro, desfez o \emph{cluster} entre duas etapas da sequência de verificações. O histórico que uma das repetições de partição pretendia preservar se perdeu, e a repetição foi redefinida. Outra queda, em 28 de outubro, interrompeu a primeira tentativa dessa repetição. As regras pré-registradas da nova medição tratam execuções interrompidas, que são descartadas e repostas, mas não compensam o atraso que novas quedas imporiam ao cronograma.\revfim{CON-09}

\section{Reflexões sobre as atividades do Trabalho de Conclusão de Curso I}
\label{sec-cronograma}

O cronograma apresentado (Tabela \ref{tab-cronograma-tcc1}), referente à disciplina de Trabalho de Conclusão de Curso~I, foi planejado com espaço de tempo de 12 semanas, organizado em seis etapas quinzenais.

\begin{table}[!h]
\captionsetup{type=table, name=Tabela}
\caption{Cronograma do Trabalho de Conclusão de Curso~I}
\centering
\small
\begin{tabular}{p{8cm}cccccc}
\hline
\textbf{Tópico} \hspace{5cm}\textbf{Semanas:} & \textbf{1--2} & \textbf{3--4} & \textbf{5--6} & \textbf{7--8} & \textbf{9--10} & \textbf{11--12} \\
\hline
Definir pergunta de pesquisa e escopo & X &  &  &  &  &  \\
Protocolo do mapeamento sistemático (bases, strings, critérios) & X &  &  &  &  &  \\
Busca e triagem (título/resumo) &  & X &  &  &  &  \\
Seleção por leitura completa, extração e classificação &  &  & X &  &  &  \\
Síntese do mapeamento e recorte do protótipo &  &  & X &  &  &  \\
Preparar ambiente experimental com Infinispan e observabilidade mínima &  &  &  & X &  &  \\
Implementar programa de carga e checagens de invariantes &  &  &  & X &  &  \\
Modelo formal com transições e canal assíncrono &  &  &  & X &  &  \\
\emph{Model checking} e análise de contraexemplos &  &  &  &  & X &  \\
\emph{Baseline} experimental com falhas selecionadas e métricas p95/p99 &  &  &  &  & X &  \\
Consolidação de resultados parciais e ajustes &  &  &  &  &  & X \\
Escrita da monografia & X & X & X & X & X & X \\
\hline
\end{tabular}
\fonte{O autor}
\label{tab-cronograma-tcc1}
\end{table}

Nas semanas 1 e 2, o objetivo era estabelecer a pergunta de pesquisa, o protocolo do mapeamento sistemático e o esquema inicial de extração e classificação. Nas semanas 3 e 4, a meta era executar a busca e a triagem inicial, e estabilizar os critérios de inclusão e exclusão. Nas semanas 5 e 6, a intenção era concluir a seleção por leitura completa, extrair dados e produzir a síntese do mapeamento, e definir o recorte implementável do protótipo. Nas semanas 7 e 8, o foco seria na preparação do ambiente experimental, com Infinispan, programa de carga e instrumentação mínima, e iniciar a especificação formal do modelo com invariantes básicas. Nas semanas 9 e 10, seria realizada a execução do \emph{model checking} inicial, sucedida por um \emph{baseline} experimental sem e com falhas selecionadas para obter dados preliminares. Nas semanas 11 e 12, previa-se consolidar resultados parciais, ajustar o recorte conforme evidências e finalizar a escrita do trabalho com metodologia, resultados parciais e planejamento das próximas etapas. A escrita do texto ocorreria em paralelo durante todo o período definido para o trabalho.



Algumas atividades desviaram do previsto e merecem comentário. O mapeamento sistemático, planejado para as semanas 5 e 6 até a extração e a classificação por leitura completa, avançou nesta etapa apenas até a seleção e a triagem: a busca correu por uma plataforma agregadora única (OpenAlex), ao passo que a consulta manual às bases originais e a classificação dos 165 estudos primários pelas dimensões D1 a D8 permaneceram para o Trabalho de Conclusão de Curso~II. A leitura das fontes também ficou parcial, com \revini{CON-19}cerca de trinta trabalhos\revfim{CON-19} lidos na íntegra; parte das fontes hospedadas atrás de autenticação foi citada a partir de resumos institucionais, e a consulta ao texto integral foi remetida à etapa seguinte. \revini{CON-24}A leitura integral concentrou-se nas referências que sustentam a fundamentação e o desenho experimental. A consulta às bases originais (ACM Digital Library, IEEE Xplore, Scopus e SpringerLink) será feita por execução manual via Portal de Periódicos CAPES.\revfim{CON-24} O ambiente experimental registrou o desvio de maior impacto operacional: o piloto de F1 correu sobre WSL2, mas o cenário F3 exigiu o módulo \texttt{sch\_netem}, ausente nesse ambiente, o que motivou a migração para uma máquina virtual Linux (Hyper-V, Ubuntu Server~24.04), decisão não antecipada no planejamento inicial.

As demais atividades cumpriram o previsto, com ressalvas de escopo. O \emph{model checking} concluiu as cinco rodadas planejadas do TLC, com contraexemplos para I1, I2, I5 e I6 e ausência de violação para I3 e I4 nas configurações exploradas. O \emph{baseline} experimental foi conduzido como piloto reduzido, limitado aos cenários de \revini{CON-22}parada ordenada de nó (F1)\revfim{CON-22} e de atraso com \emph{jitter} de canal (F3), \revini{CON-10}cujo diagnóstico consta da Seção~\ref{sec-baseline-estado}\revfim{CON-10}, e não como a bateria plena dos parâmetros T9 e T10 (30 repetições de $10^6$ operações), que segue como protocolo a escalar. Os cenários de partição (F2) e de \emph{state transfer} induzido (F4) ficaram reservados\revini{CON-27}, no planejamento do Trabalho de Conclusão de Curso~I,\revfim{CON-27} ao Trabalho de Conclusão de Curso~II. \revini{CON-11}O cronograma do Trabalho de Conclusão de Curso~I foi cumprido com duas ressalvas: a leitura completa e a classificação dos estudos do mapeamento não foram concluídas, e o \emph{baseline} das semanas 9 e 10 mostrou-se inválido como medida (Seção~\ref{sec-baseline-estado}). \revini{CON-26}No Trabalho de Conclusão de Curso~II, o conjunto de falhas foi revisto (Seção~\ref{sec-revisoes}).\revfim{CON-26}\revfim{CON-11}

\section{Próximas etapas e plano para o Trabalho de Conclusão de Curso II}
\label{sec-proximas-etapas}

O planejamento para a continuação do trabalho está organizado conforme a Tabela \ref{tab-cronograma-tcc2} e o detalhamento das atividades previstas para o Trabalho de Conclusão de Curso~II consta do Apêndice~\ref{ap-plano-tcc2}. 
\begin{table}[!h]
\captionsetup{type=table, name=Tabela}
\caption{Cronograma do Trabalho de Conclusão de Curso~II}
\centering
\small
\begin{tabular}{p{8cm}cccccc}
\hline
\textbf{Tópico} \hspace{5cm}\textbf{Semanas:} & \textbf{1--2} & \textbf{3--4} & \textbf{5--6} & \textbf{7--8} & \textbf{9--10} & \textbf{11--12} \\
\hline
Modelo formal: transições de sessão e canal assíncrono; invariantes iniciais; execução inicial de \emph{model checking} & X &  &  &  &  &  \\
Modelo formal: inclusão de falhas; refinamento de invariantes; análise de contraexemplos &  & X &  &  &  &  \\
\emph{Baseline} com falhas: cenários; cálculo de p95/p99, erro e violações &  &  & X &  &  &  \\
\emph{Baseline} com falhas: cenário de partição e, quando aplicável, rebalanceamento com \emph{state transfer}; consolidação do conjunto &  &  & X &  &  &  \\
Alinhamento formal-experimental: mapear contraexemplos para padrões observáveis; ajustes mínimos e reexecuções pontuais &  &  &  & X &  &  \\
Controlador heurístico: definição de observáveis, regras com histerese e \emph{cooldown} e perfis normal e conservador &  &  &  & X &  &  \\
Execução com controlador: repetir cenários do \emph{baseline}; coletar métricas equivalentes; registrar decisões do controlador &  &  &  &  & X &  \\
Análise comparativa: \emph{baseline} \emph{vs.} controlador e correlação com contraexemplos formais &  &  &  &  & X &  \\
Escrita da monografia & X & X & X & X & X & X \\
Consolidação final da monografia &  &  &  &  &  & X \\
\hline
\end{tabular}
\fonte{O autor}
\label{tab-cronograma-tcc2}
\end{table}

Nas semanas 1 e 2, o objetivo é consolidar o modelo formal do sistema de sessões com canal assíncrono, formalizar o conjunto inicial de invariantes e executar a primeira rodada de \emph{model checking} para obtenção de contraexemplos. Nas semanas 3 e 4, o objetivo é evoluir o modelo, incorporar falhas explícitas, refinar as invariantes e analisar os contraexemplos gerados, e delimitar quais padrões de falha devem ser reproduzidos experimentalmente. Nas semanas 5 e 6, o objetivo é executar o \emph{baseline} experimental sob falhas selecionadas (parada silenciosa de nó com reentrada, latência e \emph{jitter}, partição e, quando aplicável, rebalanceamento com \emph{state transfer}), calcular métricas como p95 e p99, taxa de erro, tempo de recuperação e taxa de violações, e consolidar o conjunto de referência. Nas semanas 7 e 8, o objetivo é alinhar modelo formal e experimento, mapear contraexemplos para padrões observáveis, realizar ajustes mínimos no programa de carga e na instrumentação para comparabilidade, e implementar o controlador heurístico integrado ao ambiente. Nas semanas 9 e 10, o objetivo é executar os mesmos cenários do \emph{baseline} com o controlador, coletar métricas equivalentes e realizar a análise comparativa, com correlação entre os resultados experimentais e os contraexemplos formais. Nas semanas 11 e 12, o objetivo é consolidar a avaliação crítica, sistematizar limitações e conclusões, e finalizar a escrita do trabalho com descrição do método, resultados, discussão e preparação para entrega e defesa. A escrita do texto ocorre em paralelo durante todo o período definido para o trabalho.

\revini{CON-12}O cronograma da Tabela~\ref{tab-cronograma-tcc2} foi modificado durante a execução. No Trabalho de Conclusão de Curso~II, as etapas das semanas 1 a 4 foram cumpridas, com a inclusão de falhas no modelo. A etapa das semanas 5 e 6 deu lugar à correção da carga e ao pré-registro da nova medição de referência, ainda não executada. O controlador, previsto para as semanas 7 e 8, foi projetado e implementado em parte antes da medição de referência, para que seu desenvolvimento não dependesse da máquina virtual. As etapas restantes, sem datas fixas, estão no Apêndice~\ref{ap-plano-tcc2}.\revfim{CON-12}

\revini{CON-31}Em síntese, a primeira evolução prevista era a inclusão de falhas explícitas no modelo formal, com análise dos contraexemplos gerados e mapeamento desses contraexemplos para padrões observáveis nos experimentos. A segunda evolução era a expansão do \emph{baseline} experimental para os cenários de partição e \emph{state transfer} induzido, e a coleta do conjunto de referência completo. A terceira evolução era o desenho e a implementação do controlador heurístico, com regras determinísticas com histerese e intervalo mínimo entre acionamentos, atuando sobre um conjunto finito de perfis pré-definidos. A última etapa era a análise comparativa entre o \emph{baseline} e a configuração com o controlador, com correlação entre os contraexemplos formais e os padrões empíricos. Desse plano, as falhas já estão no modelo formal e o controlador está projetado e implementado em parte. A segunda evolução foi modificada: o \emph{state transfer} não terá injetor próprio, e a partição e a degradação parcial serão exercidas sob carga na campanha com o cenário S3 (Seção~\ref{sec-revisoes}). O parágrafo seguinte reúne o que falta, inclusive a análise comparativa.\revfim{CON-31}

\revini{CON-13}As etapas restantes, detalhadas no Apêndice~\ref{ap-plano-tcc2}, seguem uma ordem de dependência. Na máquina virtual, falta concluir as verificações dos injetores, estabilizar o sistema operacional contra atualizações automáticas, ensaiar o F1 silencioso e recriar o \emph{cluster} com os quatro \emph{caches}. Em seguida, vêm o piloto e a campanha da nova medição de referência, que produzirão as primeiras tabelas de latência e de violação válidas, nos dois modos. Depois, a campanha com o cenário S3 exerce F2 e F4 sob carga. Em paralelo, o controlador precisa ser ligado ao laço de carga, calibrado e avaliado conforme o desenho da Seção~\ref{sec-avaliacao-controlador}, e a segunda etapa, de atuação sobre a configuração do \emph{cluster}, começa pela verificação de quais parâmetros podem mudar em tempo de execução. Na vertente formal, faltam a correção de I5 e I6 e o alinhamento entre os contraexemplos e os padrões observados nos experimentos. No mapeamento, faltam a classificação dos estudos e o fluxograma definitivo.\revfim{CON-13}
